\documentclass[11pt,a4paper]{article}
\usepackage[T1]{fontenc}
\usepackage[utf8]{inputenc}
\usepackage{lmodern}
\usepackage[margin=27mm]{geometry}
\usepackage{amsmath,amssymb,amsthm,mathtools}
\usepackage{booktabs,array,longtable}
\usepackage{microtype}
\usepackage{xcolor}
\usepackage[colorlinks=true,linkcolor=blue!45!black,citecolor=blue!45!black,urlcolor=blue!45!black]{hyperref}
\usepackage{bookmark}
\numberwithin{equation}{section}
\allowdisplaybreaks
\setlength{\parskip}{3pt}
\setlength{\parindent}{0pt}
\setlength{\emergencystretch}{2em}
\newtheorem{theorem}{Theorem}[section]
\newtheorem{proposition}[theorem]{Proposition}
\newtheorem{lemma}[theorem]{Lemma}
\newtheorem{corollary}[theorem]{Corollary}
\theoremstyle{definition}
\newtheorem{definition}[theorem]{Definition}
\theoremstyle{remark}
\newtheorem{remark}[theorem]{Remark}
\newcommand{\g}{\mathfrak g}
\newcommand{\fsl}{\mathfrak{sl}}
\newcommand{\fgl}{\mathfrak{gl}}
\newcommand{\fso}{\mathfrak{so}}
\newcommand{\fsp}{\mathfrak{sp}}
\newcommand{\tr}{\operatorname{tr}}
\newcommand{\Tr}{\operatorname{Tr}}
\newcommand{\alt}{\operatorname{Alt}}
\newcommand{\sgn}{\operatorname{sgn}}
\newcommand{\Aut}{\operatorname{Aut}}
\newcommand{\Pf}{\operatorname{Pf}}
\newcommand{\dbar}{\bar\partial}
\newcommand{\zb}{\bar z}
\newcommand{\Vht}{V_{\mathrm{ht}}}
\newcommand{\Veta}{V^\eta}
\newcommand{\Cg}{C^0_{\g}}
\newcommand{\uf}{[p_0p_1]^2f}
\newcommand{\Dthree}{[p_0p_1]^3\tr(X_0\{X_1,X_2\})}
\newcommand{\repo}{\url{https://github.com/lbotao399-stack/holomorphic-topological-feynman-diagrams}}
\title{Anomalies and All-Order Quantization of\\
Five-Dimensional Chern--Simons Theory}
\author{}
\date{November 25, 2026}
\begin{document}
\maketitle
\begin{abstract}
We study the commutative Chern--Simons theory on
$\mathbb R\times\mathbb C^2$ in the space of translation-invariant
local interactions with holomorphic derivatives and the full exterior
product of superfields. Its two-loop anomaly is evaluated by the two
theta graphs and is nontrivial for every complex simple gauge algebra
other than $\mathfrak{sl}_2$. For $\mathfrak{sl}_N$, $N\ge3$, a forced
order-one symmetric cubic cancels this anomaly. We prove that the remaining
equivariant obstruction groups vanish in the formal trace algebra,
construct the resulting local quantum master equation solution to
every order, and specialize the entire identity to every $N\ge2$.
An analytic calculation of invariant cubic tensors removes the rank
cutoffs in the no-go statement. Consequently a complex semisimple
gauge algebra admits a quantization in this sector exactly when all
its simple ideals are of type $A$.
We give the counterterms through order four, including the
transcendental fourth-order coefficient, and a complete two-loop
cancellation with an additional abelian field for
$\mathfrak{so}_5$ and other explicitly verified examples.
The all-order assertion concerns formal local perturbation theory;
the central extensions are proved through order two.
\end{abstract}
\tableofcontents

\section{Introduction and scope}

The theory considered here has a commutative algebra of spacetime
functions at classical order. Its possible quantum vertices are
determined by the quantum master equation. This distinguishes the
problem from quantizing a theory whose classical multiplication has
already been deformed by a nonzero Moyal parameter
\cite{CostelloOmega}. The analytic starting point is the decorated-graph
locality theorem of Wang and Williams \cite{WW}. The present paper
combines that input with explicit anomaly cochains and their polynomial
cohomology.

\begin{theorem}[Semisimple classification]\label{thm:main}
Let $\g$ be a complex semisimple Lie algebra equipped with a
nondegenerate invariant symmetric form. Work in the local sector
specified in Definition~\ref{def:quantization}, and let
$\epsilon,U\in\{1,-1\}$ denote the tree and loop dictionary signs
defined in Section~\ref{sec:setting}. An equivariant formal
quantization exists if and only if
\[
 \g=\bigoplus_i\mathfrak{sl}_{N_i},\qquad N_i\ge2.
\]
For one factor with the defining trace form and $N\ge3$, there are two
possible leading invariant cubic coefficients $a$, determined by
\[
 a^2=\frac{\epsilon U N^2}{256\pi^4}.
\]
Each determines a unique sequence of invariant cubic counterterms in
the formal trace algebra. The complete local identities specialize to
every $N\ge2$; this does not assert that the intrinsic finite-rank
cohomology is the stable cohomology.
\end{theorem}

The signs $\epsilon$ and $U$ record the tree and loop dictionaries and
are defined in Section~\ref{sec:setting}. Existence over $\mathbb C$ and
the classification do not depend on their choice. All numerical
coefficients displayed without these signs use the common Wick
convention $\epsilon U=1$.

The proof has three distinct parts. The two-loop calculation gives a
nonzero obstruction class. The classification of invariant cubic
tensors decides whether the quadratic first-order term can cancel
it. Finally, the complete higher obstruction calculation for type
$A$ proves that this initial cancellation extends to all orders.
The semisimple statement follows by treating the simple ideals and
the full contracted bracket, as explained in
Section~\ref{sec:lie-classification}.

\paragraph{The role of a center.}
The semisimple hypothesis is essential. Section~\ref{sec:central-extension}
constructs a mixed cubic on $\mathfrak{so}_5\oplus\mathbb C$ that
cancels the complete order-two obstruction, including every
distribution of central external inputs. This result does not
claim an all-order quantization of that central extension.

\paragraph{Relation to the computational companion.}
The original manuscript through order $\hbar^3$ \cite{RepoOriginal} contains
the full two-loop color calculations and independent certificates.
The higher-order notes \cite{RepoHigher} supply the all-order
cohomological argument and the integrated fourth-order kernels.
Here those results are organized into one theorem and supplemented
by uniform Lie-theoretic proofs, finite-rank specialization at $N=2$,
and the central-field calculation. Appendix~\ref{sec:verification}
distinguishes calculations rerun for this version from the inherited
integral data.

\section{The theory and its local cochain complex}\label{sec:setting}

Write $(\tau,z_1,z_2)$ for the coordinates of
$\mathbb R\times\mathbb C^2$ and
$\Omega=dz_1\wedge dz_2$. The superfield and its free differential are
\[
 A\in\Omega^\bullet(\mathbb R)\otimes
          \Omega^{0,\bullet}(\mathbb C^2)\otimes\g[1],
 \qquad
 DA=d\tau\,\partial_\tau A+\bar\partial A .
\]
The shift $[1]$ makes the total superfield odd. If $\kappa$ is the
invariant form, the action and the odd pairing are
\begin{equation}\label{eq:classical}
 \begin{split}
 S_{\mathrm{cl}}(A)
 &=\int\Omega\wedge
   \left(\frac12\kappa(A,DA)+\frac16\kappa(A,[A,A])\right),\\
 \omega(A,B)&=\int\Omega\wedge\kappa(A,B).
 \end{split}
\end{equation}
The bracket in the cubic combines the Lie bracket and exterior
multiplication. The pairing has degree $-1$.

\subsection{Locality and the sector being quantized}

Let $\Vht$ consist of constant-coefficient local functionals obtained
by applying finitely many holomorphic derivatives to each
superfield, taking the ordinary exterior product, and integrating
its degree-three component against $\Omega$. Functionals are taken
modulo constant-coefficient total derivatives. In particular no
real or antiholomorphic field derivatives, nor contractions of
external form indices, enter this space.

\begin{definition}\label{def:quantization}
A local formal quantization is an interaction
\[
 I+\sum_{k\ge1}\hbar^k I_k,\qquad
 I=\frac16\int\Omega\wedge\kappa(A,[A,A]),\qquad
 I_k\in\Vht,\quad \operatorname{gh}(I_k)=0,
\]
whose heat-kernel effective interactions solve the quantum master
equation at every positive scale, as formal power series.
Only positive field arity is at issue.
Equivariance means that $S/\hbar$ is invariant under
\[
 A\longmapsto\phi_g^*A,\qquad
 \hbar\longmapsto(\det g)^{-1}\hbar,\qquad g\in U(2),
\]
where $\phi_g$ acts on the two holomorphic coordinates.
\end{definition}

All superfield factors have degree one and the integral lowers
degree by three. Consequently a ghost-number-zero element of
$\Vht$ is cubic, and a ghost-number-one local obstruction is
quartic. Holomorphic derivatives do not change these degrees.
The equivariance condition on the order-$k$ cubic is
\[
 I_k(\phi_g^*A)=(\det g)^{k-1}I_k(A).
\]
For polynomial symbols this is equivalent to holomorphic
bidegree $(k,k)$ and invariance under $SL_2(\mathbb C)$.

\subsection{Momentum-conserving Taylor cochains}

The existing notation $C^0_{\g}$ denotes finite-jet
Chevalley--Eilenberg cochains of $\g[[z_1,z_2]]$, modulo the two
holomorphic total derivatives. An $r$-cochain is an alternating
multilinear function of $X_i e^{p_i\cdot z}$, polynomial in the
two-component holomorphic momenta, with
\[
 \sum_{i=0}^{r-1}p_i=0.
\]
This equation is the exact Fourier description of integration
by parts, not a restriction on a larger physical state space.
We use
\[
 [pq]=p_1q_2-p_2q_1,\qquad
 [Xe^{p\cdot z},Ye^{q\cdot z}]=[X,Y]e^{(p+q)\cdot z}.
\]
Here subscripts on a single momentum are coordinate indices.
The cochain differential, denoted by $d$, is
\begin{equation}\label{eq:CE}
 (d w)(x_0,\ldots,x_r)=
 \sum_{i<j}(-1)^{i+j}
 w([x_i,x_j],x_0,\ldots,\widehat{x_i},\ldots,
                      \widehat{x_j},\ldots,x_r).
\end{equation}
The hats omit those two arguments and $d^2=0$.
Full alternation means
$\alt w=\sum_{\sigma}\sgn(\sigma)\,w(x_{\sigma(0)},\ldots)$,
without division by a factorial.

An alternating cochain $w$ of arity $r$ gives the functional
$F_w=\frac1{r!}\int\Omega\wedge w(A,\ldots,A)$, with the
derivatives prescribed by its momentum symbol. Thus
\[
 C^{0,r}_{\g}\longrightarrow\Vht^{\,r-3}
\]
is an isomorphism of the indicated local complexes, with the
overall sign fixed below. This follows directly by polarization:
odd superfields enforce alternation, and total derivatives
enforce the sum of momenta. This dictionary concerns precisely
$\Vht$; it does not identify all translation-invariant local
functionals with this one summand.

The cochain bracket contracts one input of each cochain with
$\kappa^{-1}$ and sums the signed shuffles of the remaining
inputs. The contracted momenta are opposite, fixed by
conservation. If the two arities are $m$ and $n$, then
\begin{equation}\label{eq:cochain-bracket}
 \{a,b\}
 =\frac1{(m-1)!(n-1)!}
  \alt\left(
    a(x_1,\ldots,x_{m-1},X_e)
    \kappa^{ee'}
    b(X_{e'},x_m,\ldots,x_{m+n-2})
  \right).
\end{equation}
Here $X_e$ runs over any basis and $\kappa^{ee'}$ is the inverse
matrix of $\kappa$; the resulting contraction is basis independent.
For the classical cubic $s_0(X,Y,Z)=\kappa(X,[Y,Z])$,
\[
 \{s_0,w\}=-dw.
\]
The sign $\epsilon\in\{1,-1\}$ records the correspondence
$\{F_a,F_b\}=\epsilon F_{\{a,b\}}$. The sign $U\in\{1,-1\}$
records the common Wick sign multiplying the two-loop heat
integrals. We retain both in the existence argument.

For $J_1=F_j$ and $I_2=F_{i_2}$, the first two local equations are
\begin{equation}\label{eq:first-two-equations}
 dj=0,\qquad
 \epsilon\left(-d i_2+\frac12\{j,j\}\right)+o_2=0.
\end{equation}
Here $o_2$ is the two-loop cochain with its factor $\hbar^2$
removed. The one-loop local term vanishes, including one-loop
graphs containing holomorphically differentiated vertices.

\subsection{Heat normalization}

For the flat metric, use
\[
 D^*=-\iota_{\partial_\tau}\partial_\tau
      -4\sum_{i=1}^2\iota_{\partial_{\bar z_i}}\partial_{z_i}.
\]
The heat kernel normalized relative to $\Omega$ is
\begin{equation}\label{eq:heat}
 [K_t]_\Omega=
 \frac{\exp(-(\tau^2+|z_1|^2+|z_2|^2)/(4t))}
      {4(4\pi t)^{5/2}}\,
 d\tau\wedge d\bar z_1\wedge d\bar z_2,
 \qquad \int\Omega\wedge[K_t]_\Omega=1.
\end{equation}
The regularized propagator is
$P_\varepsilon^L=\int_\varepsilon^L D^*K_t\,dt$.
The common normalized scalar factor of a two-loop contact is
$-1/(64\pi^5)$.

Under the simultaneous dilation of coordinates by $r>0$ and
heat times by $r^2$,
\[
 r^*[K_{r^2t}]_\Omega=r^{-2}[K_t]_\Omega,\qquad
 r^*P_{r^2\varepsilon}^{r^2L}
 =r^{-2}P_\varepsilon^L.
\]
These are finite-cutoff identities. Their use for a limiting
local contact is justified by continuity of the decorated
graph distributions, not by replacing the regulated kernel
with a homogeneous expression prematurely.


\section{The exact two-loop anomaly}\label{sec:two-loop}

\subsection{The two geometric integrals}

A connected cubic graph with geometric loop number $n$ contributes a
local contact only if
\[
 2|E|=3|V|-4,\qquad |E|=3n+1,\qquad |V|=2n+2 .
\]
It then has four external fields. At two loops the two surviving
graphs are theta graphs whose three paths have lengths $(1,3,3)$
and $(2,2,3)$. The four external fields occupy their bivalent
vertices. Each graph has automorphism order four.

For a theta graph with path lengths $n_1,n_2,n_3$, the integration
over its path Schwinger parameters gives, in the present dimension,
the following expression in the Euler gamma function $\Gamma$:
\[
 \frac12
 \frac{\Gamma(\frac72-n_1)\Gamma(\frac72-n_2)
       \Gamma(\frac72-n_3)}
      {\Gamma(\frac72)(n_1-1)!(n_2-1)!(n_3-1)!}.
\]
The oriented momentum determinant supplies the graph-dependent sign.
In the orientation
of the original graph tensors the complete results are
\begin{align}
 K_{133}
 &=-\frac12\frac{\Gamma(\frac52)\Gamma(\frac12)^2}
                    {4\Gamma(\frac72)}
       [p_0p_1][p_2p_3]
   =-\frac{\pi}{20}[p_0p_1][p_2p_3],\notag\\
 K_{223}
 &=\frac12\frac{\Gamma(\frac32)^2\Gamma(\frac12)}
                    {2\Gamma(\frac72)}
       [p_0p_1][p_2p_3]
   =\frac{\pi}{30}[p_0p_1][p_2p_3].
 \label{eq:theta-kernels}
\end{align}
These are projective Schwinger integrals at zero external
displacement, with arbitrary conserved holomorphic momenta.
Their derivation by successive path integration and their original
orientation certificates are given in \cite{RepoOriginal}.

Let $W_{133}$ and $W_{223}$ denote the color contractions of the
six classical cubic tensors at the vertices, with $\kappa^{-1}$
on the seven edges, including the external-field ordering sign.
The subscript records the three path lengths. Define their full
alternating Taylor cochains by
\begin{equation}
 F_{133}=\alt\bigl([p_0p_1][p_2p_3]W_{133}\bigr),\qquad
 F_{223}=\alt\bigl([p_0p_1][p_2p_3]W_{223}\bigr).
 \label{eq:theta-cochains}
\end{equation}
With $\nu=-U/(64\pi^5)$, the automorphism factors and
\eqref{eq:theta-kernels} give
\begin{equation}
 \boxed{\quad
 o_2=\frac{\pi\nu}{240}(-3F_{133}+2F_{223})
     =-\frac{U}{15360\pi^4}(-3F_{133}+2F_{223}).
 \quad}
 \label{eq:o2}
\end{equation}
In particular the relative coefficients are $-3$ and $2$;
the scalar $U$ is common to the two graphs.

\subsection{Full color tensors for traceless matrices}

For $\g=\fsl_N$ with $\kappa=\tr$, the original contractions have
the following complete expressions:
\begin{align}
 W_{133}
 &=2\tr(X_0X_1X_2X_3)
   +(2+N^2)\tr(X_0X_1X_3X_2)\notag\\
 &\quad-4\tr(X_0X_2X_1X_3)
   +(2+N^2)\tr(X_0X_2X_3X_1)\notag\\
 &\quad-4\tr(X_0X_3X_1X_2)
   +2\tr(X_0X_3X_2X_1)
   +6N\tr(X_0X_1)\tr(X_2X_3),\notag\\
 W_{223}
 &=-2\tr(X_0X_1X_2X_3)-2\tr(X_0X_1X_3X_2)\notag\\
 &\quad+4\tr(X_0X_2X_1X_3)-2\tr(X_0X_2X_3X_1)\notag\\
 &\quad+4\tr(X_0X_3X_1X_2)-2\tr(X_0X_3X_2X_1)
       -4N\tr(X_0X_1)\tr(X_2X_3)\notag\\
 &\quad+2N\tr(X_0X_2)\tr(X_1X_3)
       +2N\tr(X_0X_3)\tr(X_1X_2).
 \label{eq:complete-theta-colors}
\end{align}
No choice of external matrices or momenta is imposed here.
Write
\[
 f(X_0,X_1,X_2)=\kappa(X_0,[X_1,X_2]),\qquad
 \beta=d\bigl([p_0p_1]^2f\bigr).
\]
In the trace normalization, full alternation of
\eqref{eq:complete-theta-colors} gives the cochain identity
\begin{equation}
 F_{133}+F_{223}=2N^2\beta,\qquad
 o_2=\pi\nu\left(\frac{F_{223}}{48}
                         -\frac{N^2\beta}{40}\right).
 \label{eq:sl-theta-boundary}
\end{equation}
Thus the bare obstruction class is
\[
 [o_2]=-\frac{U}{3072\pi^4}[F_{223}].
\]
Classical closedness follows either directly from these graph
cochains or from the first unsolved coefficient of the
regularized master equation, since the one-loop contact vanishes.

\subsection{A detector valid for every simple Lie algebra}

Use the Killing form in this subsection. For an invariant
four-cochain of momentum degree four, take the coefficient of
$[p_0p_1][p_0p_2]$. Denote its full color contraction with
$\kappa(X_0,X_1)\kappa(X_2,X_3)$ by $\ell_{kk}$, and its
contraction with $\kappa([X_0,X_1],[X_2,X_3])$ by $\ell_{ff}$.
Every contraction uses the inverse Killing form. The three
determinants $[p_0p_1],[p_0p_2],[p_1p_2]$ are algebraically
independent, so this coefficient is unambiguous.

The boundary $\beta$ has
\[
 \ell_{kk}(\beta)=2\dim\g,\qquad
 \ell_{ff}(\beta)=\dim\g.
\]
Consequently the linear functional
\[
 \varpi=(\dim\g)\ell_{kk}-2(\dim\g)\ell_{ff}
 \qquad\text{satisfies}\qquad \varpi(\beta)=0 .
\]
It annihilates every boundary in the invariant sector in question.
Indeed exactness of invariants for the finite-dimensional
representations of the adjoint group and $SL_2$ at fixed momentum degree supplies
an invariant primitive whenever an invariant cochain is exact.
A cubic primitive of degree four has momentum polynomial
$[p_0p_1]^2$ and an invariant alternating color tensor.
That tensor is a multiple of $f$ for every simple algebra.
Thus the entire invariant boundary space is $\mathbb C\beta$.

To evaluate the detector, let $e_b,e^b$ be Killing-dual bases,
let $P$ exchange the two factors of $\g\otimes\g$, and put
\[
 \Omega_{\mathrm{ad}}
   =\sum_b\operatorname{ad}(e_b)\otimes\operatorname{ad}(e^b),
 \qquad n_4=\Tr_{\g\otimes\g}(P\Omega_{\mathrm{ad}}^4).
\]
Contraction of the original theta networks gives
\begin{align}
 \ell_{kk}(F_{223})&=-6\dim\g,&
 \ell_{ff}(F_{223})&=-8n_4,\notag\\
 \varpi(F_{223})&=2\dim\g(8n_4-3\dim\g),&
 \varpi(F_{133})&=-\varpi(F_{223}).
 \label{eq:universal-detector}
\end{align}
For clarity, these contractions can be reduced to the three
pairings $(01)(23),(02)(13),(03)(12)$ before alternation.
For $W_{223}$ their metric contractions are
$(-\dim\g/2,\dim\g/4,\dim\g/4)$ and their two-bracket contractions
are $(0,n_4,n_4)$. Full alternation applies the linear combinations
$(8,-4,-4)$ and $(0,-4,-4)$, respectively, giving
\eqref{eq:universal-detector}.

Here is the exact family-independent evaluation of the remaining
Casimir trace. In Vogel's normalization let
$(v_1,v_2,v_3)$ be the three parameters and set
$t_V=v_1+v_2+v_3$. The fourth adjoint Casimir formula
\cite{MSV} gives
\begin{align}
 n_4&=\dim\g\left[
  \frac{3t_V(v_1^2+v_2^2+v_3^2)
              -(v_1^3+v_2^3+v_3^3)}
       {16t_V^3}-\frac18\right],\notag\\
 8n_4-3\dim\g
 &=-\frac{3\dim\g\,(t_V+v_1)(t_V+v_2)(t_V+v_3)}
             {2t_V^3}.
 \label{eq:vogel-detector}
\end{align}
The substitutions and results are displayed together:
\begin{center}
\renewcommand{\arraystretch}{1.55}
\begin{tabular}{@{}lll@{}}
\toprule
$\g$ & $(v_1,v_2,v_3)$ & $8n_4-3\dim\g$\\
\midrule
$\fsl_N$ & $(-2,2,N)$
  & $-\dfrac{3(N^2-1)(N^2-4)}{N^2}$\\
\addlinespace[3pt]
$\fso_N$ & $(-2,4,N-4)$
  & $-\dfrac{3N(N-1)(N-3)(N-4)(N+2)}{2(N-2)^3}$\\
\addlinespace[3pt]
$\fsp_N$ & $(-2,1,\frac{N+4}{2})$
  & $-\dfrac{3N(N-2)(N+1)(N+3)(N+4)}{2(N+2)^3}$\\
\addlinespace[3pt]
$G_2$ & $(-2,\frac{10}{3},\frac83)$ & $-\dfrac{385}{12}$\\
\addlinespace[3pt]
$F_4$ & $(-2,5,6)$ & $-\dfrac{12740}{81}$\\
\addlinespace[3pt]
$E_6$ & $(-2,6,8)$ & $-\dfrac{975}{4}$\\
\addlinespace[3pt]
$E_7$ & $(-2,8,12)$ & $-\dfrac{34580}{81}$\\
\addlinespace[3pt]
$E_8$ & $(-2,12,20)$ & $-\dfrac{12152}{15}$\\
\bottomrule
\end{tabular}
\end{center}
For $\fsp_N$ the defining representation has even dimension $N$.
The zeros in the simple families occur precisely at
$\fsl_2=\fso_3=\fsp_2$; $\fso_4$ is not simple.

\begin{theorem}[Nontriviality of the bare anomaly]
\label{thm:bare-anomaly}
The bare two-loop anomaly is a nonzero cohomology class for every
complex simple gauge algebra other than $\fsl_2$.
\end{theorem}
\begin{proof}
Equations \eqref{eq:o2} and \eqref{eq:universal-detector} give
\[
 \varpi(o_2)=\frac{\pi\nu}{24}
                 \dim\g(8n_4-3\dim\g).
\]
The table makes this nonzero in the stated cases, whereas
$\varpi$ vanishes on every possible invariant boundary.
In Killing normalization for $\fsl_2$, direct contraction gives
$F_{133}=\frac72\beta$ and $F_{223}=-3\beta$; hence its
obstruction is exact. Rescaling the invariant form by $c\ne0$
multiplies each six-vertex, seven-edge theta graph by
$c^{6-7}=c^{-1}$ and cannot change nontriviality.
\end{proof}

\section{The gauge-algebra restriction at every rank}
\label{sec:lie-classification}

The restriction on the simple gauge algebra is uniform in its rank.  The
representation-theoretic input below is classical; equation~(18) of
\cite{UrichukWalton} records its more general highest-weight version.  We
give the argument needed here and then determine the entire first-order
cohomology in the absence of a symmetric cubic invariant.

\begin{proposition}[Invariant cubic tensors]
\label{prop:uniform-cubic}
For a complex simple Lie algebra,
\begin{equation}
\dim(\g^{*\otimes3})^{\g}=
\begin{cases}
2,&\g=\mathfrak{sl}_N,\quad N\ge3,\\
1,&\text{every other simple }\g.
\end{cases}
\label{eq:uniform-cubic}
\end{equation}
In every case the alternating part is spanned by
$f(X,Y,Z)=\kappa(X,[Y,Z])$.  For $\mathfrak{sl}_N$, $N\ge3$, the additional
tensor is the symmetric tensor $\tr(X\{Y,Z\})$.
\end{proposition}

\begin{proof}
Let $\theta$ be the highest root, $\rho$ the half-sum of positive roots,
and $W$ the Weyl group.  Write $m_{\g}$ for the weight multiplicity in
the adjoint representation.  The invariant form identifies the adjoint
representation with its dual.  The Weyl character formula therefore gives
\begin{align}
\dim(\g^{*\otimes3})^{\g}
&=\dim\operatorname{Hom}_{\g}(\g,\g\otimes\g)\notag\\
&=\sum_{w\in W}\det(w)\,
m_{\g}\bigl(\theta+\rho-w(\theta+\rho)\bigr).
\label{eq:adjoint-weyl-sum}
\end{align}
Here $m_{\g}(0)=\operatorname{rank}\g$, every root has multiplicity one,
and all other weights have multiplicity zero.  Since $\theta+\rho$ lies
strictly inside the dominant chamber, only $w=1$ gives the zero weight.

Suppose that $\alpha=\theta+\rho-w(\theta+\rho)$ is a root.
Dominance implies that $\alpha$ is positive.  Preservation of the inner
product by $W$ implies
\begin{equation}
2(\theta+\rho,\alpha)=(\alpha,\alpha),\qquad
\langle\theta+\rho,\alpha^\vee\rangle=1.
\end{equation}
The positive coroot $\alpha^\vee$ has positive integral coefficients in
the basis of simple coroots.  Each such coefficient contributes at least
one to its pairing with $\rho$, and its pairing with the dominant
weight $\theta$ is nonnegative.  Consequently $\alpha$ is a simple root
with $\langle\theta,\alpha^\vee\rangle=0$.  Conversely each such simple
root contributes, with $w=s_\alpha$ and sign $-1$.  Regularity of
$\theta+\rho$ ensures that no second Weyl element gives the same image.
Thus
\begin{equation}
\dim(\g^{*\otimes3})^{\g}
=\operatorname{rank}\g-
\#\{i:\langle\theta,\alpha_i^\vee\rangle=0\}
=\#\{i:\langle\theta,\alpha_i^\vee\rangle\ne0\}.
\label{eq:highest-root-support}
\end{equation}
In type $A_r$ with $r\ge2$ the highest root is the sum of the first and
last fundamental weights.  In type $A_1$ and type $C_r$ it is twice one
fundamental weight.  In types $B_r$ with $r\ge3$, $D_r$ with $r\ge4$,
and all exceptional types it is one fundamental weight; $B_2=C_2$ is
already covered.  This proves~\eqref{eq:uniform-cubic}.
The tensor $f$ is nonzero and alternating.  In type $A_r$, $r\ge2$,
$\tr(X\{Y,Z\})$ is nonzero and symmetric, so these two tensors exhaust
the invariant space.  Outside this family $f$ exhausts it.
\end{proof}

\begin{lemma}[Polynomial cocycles in the alternating color channel]
\label{lem:f-channel-all-degrees}
Suppose $\g$ is simple and $(\g^{*\otimes3})^{\g}=\mathbb C f$.
In the conserved-momentum complex $\Cg$,
\begin{equation}
H^3(\Cg)_{(0,0)}=\mathbb C[s_0],\qquad
H^3(\Cg)_{(a,b)}=0\quad\text{for }a+b>0.
\label{eq:f-channel-H3}
\end{equation}
The second statement holds at every polynomial derivative order, without
rotation or dilation equivariance.
\end{lemma}

\begin{proof}
At fixed momentum degree the cochain spaces are finite-dimensional
representations of the adjoint group.  The action of a constant
$X\in\g$ is $d\iota_X+\iota_Xd$ and hence is zero on cohomology.
Complete reducibility and exactness of the invariants functor imply that
every class has an invariant representative and every invariant boundary
has an invariant primitive.  It therefore suffices to work with invariant
cochains.

Write an invariant three-cochain as
\begin{equation}
P(p,q)\,\kappa(X,[Y,Z]),\qquad p+q+r=0.
\end{equation}
Alternation of the complete cochain is equivalent to
\begin{equation}
P(p,q)=P(q,p)=P(p,-p-q).
\label{eq:P-permutations}
\end{equation}
The three four-color contractions
\[
\begin{gathered}
\kappa([X_1,X_2],[X_3,X_4]),\qquad
\kappa([X_1,X_3],[X_2,X_4]),\\
\kappa([X_1,X_4],[X_2,X_3])
\end{gathered}
\]
have only the Jacobi relation.
For example, in any embedded $\mathfrak{sl}_2$ choose two elements with
$\kappa([X,Y],[X,Y])\ne0$.  The assignments $(X,X,Y,Y)$ and
$(X,Y,Y,X)$ make the first and third contractions independent.
The Chevalley--Eilenberg differential consequently vanishes exactly when
\begin{equation}
P(p,q)+P(p+q,r)=P(q,r)+P(p,q+r).
\label{eq:P-cocycle}
\end{equation}
The momenta $p,q,r$ in this identity are independent; the fourth momentum
is $-p-q-r$.

Set $q=0$ in~\eqref{eq:P-cocycle}.  Then $P(p,0)=P(0,0)$.
For a tangent vector $u$, let
$D_2P(v,0)[u]$ denote the directional derivative in the second argument.
Differentiating~\eqref{eq:P-cocycle} in $r$ at zero gives
\begin{equation}
D_2P(p,q)[u]=D_2P(p+q,0)[u]-D_2P(q,0)[u].
\end{equation}
Symmetry in $p,q$ shows that the polynomial one-form
$D_2P(v,0)$ is closed.  Its polynomial primitive, vanishing at zero, is
\begin{equation}
H(v)=\int_0^1D_2P(tv,0)[v]\,dt.
\end{equation}
The preceding derivative identity and evaluation at $q=0$ imply
\begin{equation}
P(p,q)=P(0,0)+H(p+q)-H(p)-H(q).
\label{eq:P-primitive}
\end{equation}
The second symmetry in~\eqref{eq:P-permutations} now gives
$H(v)+H(-v)=H(0)+H(0)=0$.  Hence the two-cochain
\begin{equation}
-H(p)\,\kappa(X,Y),\qquad p+q=0,
\end{equation}
is alternating.  With the differential convention used throughout,
its differential is
\begin{equation}
\bigl[H(p+q)-H(p)-H(q)\bigr]\kappa(X,[Y,Z]).
\end{equation}
It removes every positive-degree term of $P$.  The remaining constant
multiple of $s_0$ is not exact: a momentum-zero invariant two-tensor is
symmetric, and therefore there is no alternating invariant two-cochain
of that degree.  All constructions preserve the momentum bidegree.
\end{proof}

\begin{theorem}[Uniform two-loop obstruction]
\label{thm:uniform-no-go}
Let $\g$ be a complex simple Lie algebra outside the family
$\mathfrak{sl}_N$, $N\ge2$, with any nondegenerate invariant bilinear
form.  In $\Vht$ there is no solution of the quantum master equation
through order $\hbar^2$, for any choice of closed first-order local
counterterm.  Neither adjoint invariance of that counterterm nor
$U(2)$-equivariance is required.
\end{theorem}

\begin{proof}
By Proposition~\ref{prop:uniform-cubic} and
Lemma~\ref{lem:f-channel-all-degrees}, every closed first-order Taylor
cochain $j$ is cohomologous to a constant multiple of $s_0$.
The bracket induces a bracket on cohomology, and
$\{s_0,s_0\}=0$.  Therefore
\begin{equation}
\left[\frac12\{j,j\}\right]=0.
\end{equation}
The order-two master equation requires this class to cancel $[o_2]$.
The two-loop detector is nonzero for every simple algebra other than
$\mathfrak{sl}_2$: with Killing normalization its value on $F_{223}$ is
$2\dim\g\,(8n_4-3\dim\g)$, with the nonzero expressions displayed in
the two-loop computation.  Rescaling the invariant form multiplies the
obstruction by a nonzero scalar.  Hence $[o_2]\ne0$, a contradiction.
This uses the exact two-loop calculation, not a finite-rank extrapolation
of the first-order cohomology.
\end{proof}

\subsection{Semisimple direct sums}

Let $\g=\bigoplus_i\g_i$ be semisimple, with each $\g_i$ simple.
Invariance and $[\g_i,\g_i]=\g_i$ give
\begin{equation}
\kappa(\g_i,\g_j)=0\quad(i\ne j),\qquad
(\g^{*\otimes3})^{\g}
=\bigoplus_i(\g_i^{*\otimes3})^{\g_i}.
\label{eq:semisimple-splitting}
\end{equation}
For the second identity, every distribution of three tensor factors
over more than one ideal contains an ideal occurring exactly once.
Its contribution vanishes because $(\g_i^*)^{\g_i}=0$.

\begin{corollary}[Semisimple classification in the specified sector]
\label{cor:semisimple-classification}
Under the local heat-kernel hypotheses of
Theorem~\ref{thm:all-orders-traceless}, a semisimple theory admits an
all-order $U(2)$-equivariant quantization with counterterms in $\Vht$
if and only if
\begin{equation}
\g=\bigoplus_i\mathfrak{sl}_{N_i},\qquad N_i\ge2.
\end{equation}
If a non-type-$A$ simple ideal occurs, the obstruction already appears
at order $\hbar^2$, even without equivariance of the counterterms.
\end{corollary}

\begin{proof}
Choose an invariant representative of the first-order cohomology class.
Its Taylor cochain splits over the ideals by
\eqref{eq:semisimple-splitting}.  The inverse invariant form is block
diagonal, so contractions between cubic cochains supported on distinct
ideals vanish.  The bare connected two-loop graphs likewise have all
their color labels in one ideal.  Evaluating the order-two equation on
four arguments from any non-type-$A$ ideal therefore gives precisely
the impossible simple-algebra equation of
Theorem~\ref{thm:uniform-no-go}.  Evaluating a Chevalley--Eilenberg
coboundary on that ideal remains a coboundary.  This argument uses
splitting of invariant cubic tensors; it does not assume that arbitrary
restrictions commute with the Batalin--Vilkovisky bracket.

Conversely, take the quantization of each $\mathfrak{sl}_{N_i}$ supplied
by Theorem~\ref{thm:all-orders-traceless} and sum the actions.  Cross
brackets and cross propagator contractions vanish, so the full master
equation is the sum of the individual master equations.  If the invariant
form on a factor is $c\tr$ with $c\ne0$, replace its trace-normalized
quantum action $S(\hbar)$ by $cS(\hbar/c)$.  A connected graph contributes
$c$ from each vertex and $c^{-1}$ from each internal contraction.
Together with the counterterm orders, its coefficient of $\hbar^n$
is multiplied by $c^{1-n}$, so this replacement preserves every
coefficient of the master equation.  It handles every nondegenerate
invariant form on the semisimple algebra.
\end{proof}

\begin{remark}
The exclusion of an abelian center is essential to the necessity proof.
For a reductive algebra with center, a mixed tensor consisting of the
invariant bilinear form on a simple ideal and a linear functional on the
center is allowed.  Two such cubic vertices can contract through the
center and contribute to the four-simple-argument equation.  Thus
\eqref{eq:semisimple-splitting} does not extend to reductive algebras.
\end{remark}

\section{The first three quantum counterterms}
\label{sec:first-counterterms}

Throughout this section $\g=\fsl_N$ and $\kappa=\tr$.
To retain the original cochain notation, write
\[
 D(X_0e^{p_0z},X_1e^{p_1z},X_2e^{p_2z})
   =[p_0p_1]\tr(X_0\{X_1,X_2\}),\qquad
 p_0+p_1+p_2=0.
\]
Here $D$ denotes a cubic cochain, whereas the $D$ acting on
superfields in \eqref{eq:classical} is the free differential.
The matrix anticommutator is
$\{X_1,X_2\}=X_1X_2+X_2X_1$.
The color tensor is symmetric and the momentum determinant is
alternating, so the complete cochain is alternating.

\subsection{The forced order-one interaction}

At first order the only invariant equivariant cubic is $j=aD$.
Its classical differential vanishes. Contraction of the full
shuffle bracket, using the traceless completeness relation, gives
\begin{equation}
 \frac12\{D,D\}
  =\frac{F_{223}}{12N^2}
     +\frac{N^2+4}{2N^2}\,\beta .
 \label{eq:DD}
\end{equation}
Substitution in \eqref{eq:first-two-equations} yields
\begin{align}
 d i_2
 &=\left(\frac{a^2}{12N^2}
                       +\frac{\epsilon\pi\nu}{48}\right)F_{223}
   +\left(\frac{a^2(N^2+4)}{2N^2}
                       -\frac{\epsilon\pi\nu N^2}{40}\right)\beta .
 \label{eq:order-two-assembly}
\end{align}
For $N\ge3$, Theorem~\ref{thm:bare-anomaly} forces the first
coefficient to vanish. Therefore
\begin{align}
 a^2&=-\frac{\epsilon\pi\nu N^2}{4}
       =\frac{\epsilon U N^2}{256\pi^4},\notag\\
 i_2&=\frac{\epsilon U(3N^2+10)}{1280\pi^4}
         [p_0p_1]^2\tr(X_0[X_1,X_2]).
 \label{eq:first-two-counterterms}
\end{align}
These are equalities of full cochains, not only cohomology classes.
In the convention $\epsilon U=1$, the two roots are
$a=N/(16\pi^2)$ and $a=-N/(16\pi^2)$.

\subsection{The marked two-loop contribution at order three}

There are three sources at order $\hbar^3$: the classical
differential of $i_3$, the tree bracket of $j$ and $i_2$, and the
two-loop theta graphs with one $j$ vertex. Bare three-loop graphs
and one-loop graphs with any of these vertices vanish.
The complete differentiated theta calculation gives
\begin{align}
 \{D,[p_0p_1]^2f\}
  &=\frac13d\bigl([p_0p_1]^3\tr(X_0\{X_1,X_2\})\bigr),\notag\\
 \Theta(D)
  &=-\frac{N^2}{60}
       d\bigl([p_0p_1]^3\tr(X_0\{X_1,X_2\})\bigr).
 \label{eq:marked-theta}
\end{align}
Here $\Theta(D)$ is the sum over both theta graphs and all
single marked vertices, with the common scalar $\pi\nu$
removed. Its graph automorphism factors are included.
The local equation is consequently
\[
 d i_3=\{j,i_2\}+\epsilon\pi\nu\,a\,\Theta(D).
\]
Using \eqref{eq:first-two-counterterms} and \eqref{eq:marked-theta},
\begin{align}
 i_3
 &=\frac{a\epsilon U}{\pi^4}
       \left(\frac{3N^2+10}{3840}+\frac{N^2}{3840}\right)
          [p_0p_1]^3\tr(X_0\{X_1,X_2\})\notag\\
 &=\frac{a\epsilon U(2N^2+5)}{1920\pi^4}
          [p_0p_1]^3\tr(X_0\{X_1,X_2\}).
 \label{eq:i3}
\end{align}
The tensor identities in this section were independently replayed
from the original graphs using exact rational matrices; their
finite evaluations are certificates accompanying the formal
trace identities, rather than substitutes for them.

\subsection{Rank two and the two branches}

For traceless $2\times2$ matrices,
\[
 X_1X_2+X_2X_1=\tr(X_1X_2)\mathbf1_2,\qquad
 \tr(X_0\{X_1,X_2\})=0 .
\]
Evaluation of the full identities at $N=2$ therefore gives
\begin{equation}
 j=0,\qquad
 i_2=\frac{11\epsilon U}{640\pi^4}
             [p_0p_1]^2\tr(X_0[X_1,X_2]),\qquad
 i_3=0.
 \label{eq:rank-two}
\end{equation}
It is not necessary to cancel a nonzero class at rank two:
the bare anomaly there is already exact.

The Lie automorphism $X\mapsto-X^{\mathsf T}$ preserves the trace
pairing and fixes $f$, while reversing the symmetric cubic.
It commutes with the cochain differential and every graph
contraction. Once the all-order uniqueness in the formal sector
is proved, it follows that the two choices of $a$ are exchanged
by this automorphism: every even-order counterterm is fixed and
every odd-order counterterm changes sign. Thus for $N\ge3$
there are two leading couplings in the fixed color convention,
related by a symmetry of the underlying classical theory.
At $N=2$ every odd-order cubic vanishes, so the two formal
branches specialize to the same quantization.

\section{All-order quantization in the holomorphic local sector}
\label{sec:all-orders}

The order of a counterterm is its exponent of $\hbar$; it need not equal
the geometric loop number of a graph containing that counterterm.
Throughout this section the coefficients are complex, and the two sign
conventions $\epsilon,U\in\{1,-1\}$ are retained.  No choice of either
sign is needed for the existence theorem.

\subsection{The local contact stays in the required space}

The space $\Vht$ consists of integrated exterior products of holomorphic
jets of the superfield, with constant coefficients.  Integration uses the
full degree-three real and antiholomorphic exterior component.  In
particular, a ghost-number-zero element is cubic and a
ghost-number-one element is quartic.  Membership in $\Vht$ is a condition
on a local cochain, stronger than belonging to its image in cohomology.

\begin{lemma}[Closure of holomorphic local vertices under contacts]
\label{lem:full-wedge-contacts}
For the flat heat kernel on $\mathbb R\times\mathbb C^2$, every
total-collapse contact constructed from vertices in $\Vht$ belongs to
$\Vht$.  This includes vertices with arbitrary finite holomorphic
derivatives.  A contact of odd geometric loop number vanishes.
\end{lemma}

\begin{proof}
An integration by parts expresses every derivative at a vertex as a
finite sum of holomorphic derivatives on incident kernels and external
test fields.  These are the decorated graphs of
Wang--Williams~\cite[Propositions 3.18, 3.21--3.23]{WW}.  Their
compactified graph distributions are continuous on compactly supported
smooth test forms.  A nonzero total-collapse contact is supported on
the diagonal and differentiates its test form only holomorphically.

It remains to check its exterior structure.  Write $|V|$ and $|E|$ for
the numbers of vertices and internal edges of a connected collapsing
graph, and $n=|E|-|V|+1$ for its geometric loop number.  The tight
condition in three real and antiholomorphic exterior directions is
\begin{equation}
  2|E|=3|V|-4,
  \qquad |E|=3n+1,
  \qquad |V|=2n+2.
  \label{eq:contact-tight-degree}
\end{equation}
The extended kernel has exterior degree three on each edge.
After extracting the $|E|-1$ Schwinger differentials, its spatial
exterior degree is
\begin{equation}
  3|E|-(|E|-1)=2|E|+1=3(|V|-1).
  \label{eq:relative-form-saturation}
\end{equation}
All these differentials are differences between vertex coordinates.
They therefore saturate the relative real and antiholomorphic
directions.  Only the diagonal pullback of the external exterior
product remains.  No contraction of an external form index and no
additional background form is required.  The resulting local vertex
is precisely a full exterior product of holomorphic jets.  For cubic
inputs its external arity is
\begin{equation}
  3|V|-2|E|=4.
\end{equation}
Self-loops give zero: their differential factors are coordinate
differences, whose diagonal pullback remains zero after holomorphic
differentiation.  Finally, reflection of the real relative coordinates
changes the contact by $(-1)^n$.  Its local holomorphic test data are
unchanged.  This proves the odd-loop assertion.
\end{proof}

The regulated characters also record covariance and the limiting
step.  Let $r>0$ and
$|u_1|=|u_2|=1$, and pull back by
$(\tau,z_1,z_2)\mapsto(r\tau,ru_1z_1,ru_2z_2)$.  Direct substitution
in the normalized kernel gives
\begin{align}
  \phi_{r,u}^*[K_{r^2t}]_\Omega
    &=r^{-2}(u_1u_2)^{-1}[K_t]_\Omega,\notag\\
  \phi_{r,u}^*[D^*K_{r^2t}]_\Omega
    &=r^{-4}(u_1u_2)^{-1}[D^*K_t]_\Omega,\notag\\
  \phi_{r,u}^*P_{r^2\varepsilon}^{r^2L}
    &=r^{-2}(u_1u_2)^{-1}P_\varepsilon^L.
  \label{eq:regulated-covariance-all-orders}
\end{align}
These identities hold at positive cutoffs.  Continuity of pullback
on test forms passes them through the distributional boundary
restriction supplied by the compactified graph theorem.  Thus the
local obstruction has the prescribed rotation and dilation
characters.  Together with the exterior-form argument, this gives
the required statement at cochain level; exactness
in a larger space of local functionals alone would not prove it.

\subsection{The degree-four cohomology at every momentum degree}

We first work in the formal invariant trace algebra for traceless
matrices: distinct products of cyclic trace words are independent,
and a one-letter trace is zero.  For cohomology in degree four this
algebra agrees with the actual invariant tensors of $\fsl_N$ for
$N\ge5$, since only arities three, four, and five occur.  The
differential inserts a commutator and preserves the number of traces.

\begin{proposition}[The complete obstruction module]
\label{prop:all-order-h4}
For $w\ge2$, the formal trace complex satisfies
\begin{equation}
  H^3(C^0_{\fsl_N})_{\deg=2w}=0,
  \qquad
  H^4(C^0_{\fsl_N})_{\deg=2w}
       \cong\operatorname{Sym}^{2w-4}(\mathbb C^2).
  \label{eq:all-order-cohomology}
\end{equation}
The second space is represented entirely by two-trace cochains.
For the trace word $\tr(X_0X_1)\tr(X_2X_3)$, put
$P=p_0+p_1=-p_2-p_3$.  Its coefficient is
\begin{equation}
  [p_0P][p_2P]h(P),
  \qquad h\text{ homogeneous of degree }2w-4.
  \label{eq:all-order-double-trace}
\end{equation}
The other pairings are fixed by alternation.  Consequently
\begin{equation}
 \dim H^4(C^0_{\fsl_N})_{\deg=2w}=2w-3,
 \qquad
 \dim\bigl(H^4(C^0_{\fsl_N})_{\deg=2w}\bigr)^{SL_2}
 =\begin{cases}1,&w=2,\\0,&w\ge3.\end{cases}
 \label{eq:all-order-invariant-vanishing}
\end{equation}
\end{proposition}

\begin{proof}
\emph{One trace.}
A trace word of arity $k$ is a closed polygon with $k$ vertices:
write its momenta as $p_i=x_{i+1}-x_i$, with $x_k=x_0$.
The word coefficient is a polynomial in $x_0,\ldots,x_{k-1}$,
invariant under their simultaneous translation and cyclic with sign
$(-1)^{k-1}$.  Splitting one input into a commutator in the
Chevalley--Eilenberg differential adds two adjacent sides.  Equivalently,
the differential deletes the intervening polygon vertex.  Up to a
uniform degree sign this is the alternating deletion differential.

Here is the grading and acyclicity argument explicitly.  Let $j=k-1$.
Before imposing cyclicity, translation-invariant polynomials on
$j+1$ vertices have the ordinary additive bar differential on their
$j$ successive independent increments.  Its end faces omit the first
or last increment, and its interior faces add two consecutive
increments.  At fixed polynomial degree $m$, normalization removes
degenerate cochains.  The resulting complex is the graded dual of the
reduced bar complex of a polynomial algebra on two generators:
the coproduct $f(p+q)$ becomes multiplication after dualizing in the
factorial basis.  Write $HH$ for the cohomology of this full deletion
complex and $HC$ for the cohomology of its signed cyclic subcomplex.
The length-two Koszul resolution consequently gives
\begin{equation}
 HH^j_m=
 \begin{cases}
  \Lambda^j(\mathbb C^2),&m=j,\quad 0\le j\le2,\\
  0,&\text{otherwise}.
 \end{cases}
 \label{eq:ordinary-bar-bidegrees}
\end{equation}
This is a statement about translation-invariant polygon polynomials,
with no remaining total momentum variable.

The cyclic bicomplex is available degree by degree.  Rotation, each
face, and insertion of a zero side preserve $m$.  The complex with
its last face omitted is contracted by appending a duplicate final
vertex, with the alternating degree sign.  The finite cyclic norm
identifies cyclic invariants and coinvariants over characteristic
zero.  Thus its long exact sequence contains
\begin{equation}
  HC^{j-2}_m\longrightarrow HC^j_m\longrightarrow HH^j_m.
  \label{eq:polygon-cyclic-sequence}
\end{equation}
For $m>2$, the last space vanishes at every $j$ by
\eqref{eq:ordinary-bar-bidegrees}.  Starting with zero in negative
cochain degrees and inducting in $j$ proves $HC^j_m=0$ for every
$j$ and $m>2$.  The single-trace summand in Chevalley--Eilenberg
arity $k$ is $HC^{k-1}_m$, not $HC^k_m$.
Deleting the one-letter trace in the traceless complex cannot change
cohomology in arities three and four: their incoming and outgoing
maps involve only arities two through five.  Hence the one-trace
summands of both groups in \eqref{eq:all-order-cohomology} vanish.

\emph{Two traces.}
At arity four the only remaining word type is $(01)(23)$.
Write its coefficient as $F(r,q,P)$ on the momenta
$(r,P-r,q,-P-q)$.  Alternation gives
\begin{align}
 F(P-r,q,P)&=-F(r,q,P),\notag\\
 F(r,-P-q,P)&=-F(r,q,P),\notag\\
 F(q,r,-P)&=F(r,q,P).
 \label{eq:pair-exchanges}
\end{align}
The coefficient of an oriented three-letter trace times the other
two-letter trace in $dF$ says
\begin{equation}
 F(r,q,P)+F(s,q,P)+F(P-r-s,q,P)=0.
\end{equation}
Using \eqref{eq:pair-exchanges}, this is precisely
\begin{equation}
 F(r+s,q,P)=F(r,q,P)+F(s,q,P).
 \label{eq:pair-additivity}
\end{equation}
Polynomial additivity implies linearity in $r$.  The first exchange
then forces that linear functional to annihilate $P$.  The two
coordinates of $P$ are relatively prime in their polynomial ring;
therefore a polynomial linear functional annihilating $P$ is a
polynomial multiple of $[rP]$.  Applying the same argument to $q$
gives
\begin{equation}
 F(r,q,P)=[rP][qP]h(P),\qquad h(-P)=h(P).
\end{equation}
Conversely this expression satisfies every coefficient of $dF=0$.
At total momentum degree $2w$, the polynomial $h$ has degree
$2w-4$, so its evenness is automatic.  There is no arity-three
two-trace word on traceless matrices.  These cocycles are therefore
not boundaries.  Their $SL_2$ representation and dimension are those
of binary forms of degree $2w-4$.  This proves all assertions.
\end{proof}

\subsection{Induction and finite-rank specialization}

\begin{theorem}[All-order traceless quantization]
\label{thm:all-orders-traceless}
Take the defining trace form on $\fsl_N$ and the flat heat-kernel
scheme.  For each of the two complex solutions of
\begin{equation}
 a^2=\epsilon U\frac{N^2}{256\pi^4},
 \label{eq:all-order-leading-root}
\end{equation}
the first-order cubic $J_1=aF_D$ has a formal local quantization
\begin{equation}
 S_{\mathrm{cl}}+\hbar J_1+\sum_{w\ge2}\hbar^w I_w,
 \qquad I_w\in\Vht,
 \label{eq:all-order-action}
\end{equation}
with the prescribed $U(2)$ characters.  In the formal trace algebra
each $I_w$ is the unique invariant cubic solving the order-$w$
equation.  The complete identities specialize to every integer
$N\ge2$.  At $N=2$, $J_1$ is identically zero; the conclusion there
is existence and does not infer intrinsic uniqueness from the formal
trace calculation.
\end{theorem}

\begin{proof}
The two-loop calculation supplies $J_1$ and $I_2$ satisfying the local
equation through order two.  Suppose the local equations below order
$w\ge3$ have been solved.  The leading remaining obstruction is
classically closed.  One way to see this is to use the finite-scale
quantum Batalin--Vilkovisky identity.  With $Q$ the free differential,
$\Delta_L$ the heat-regularized contraction, and $S[L]$ the effective
interaction, its curvature satisfies
\begin{gather}
 O[L]=QS[L]+\hbar\Delta_LS[L]
                  +\tfrac12\{S[L],S[L]\}_L,\notag\\
 \bigl(Q+\hbar\Delta_L+\{S[L],-\}_L\bigr)O[L]=0.
 \label{eq:finite-scale-bianchi}
\end{gather}
At the first unsolved order the lower curvatures vanish.  Passing to
the local contact, or equivalently applying the local obstruction
construction of~\cite{CostelloBV}, gives classical closedness.
The compactified boundary factorization of~\cite[Proposition
3.22]{WW} identifies this local equation with the effective master
equation at positive scale.

Lemma~\ref{lem:full-wedge-contacts} puts the actual local obstruction
in $\Vht$, so its Taylor cochain is an element of
$C^{0,4}_{\fsl_N}$.  The determinant character gives total momentum
degree $2w$.  Its $SU(2)$ invariance is $SL_2$ invariance because the
momentum coefficients are holomorphic polynomials.  By
\eqref{eq:all-order-invariant-vanishing}, its invariant cohomology
class is zero.  Complete reducibility in each finite degree gives
an invariant primitive.  The chain identification with $\Vht$
turns this primitive into the required local counterterm $I_w$.

To verify uniqueness within the stated formal sector, three
conserved momenta leave two independent momentum vectors.
Their invariant polynomial ring is generated by their determinant.
The cubic Taylor cochain at order $w$ is therefore a scalar multiple
of
\begin{equation}
 [p_0p_1]^w
 \begin{cases}
  \tr\bigl(X_0[X_1,X_2]\bigr),&w\text{ even},\\
  \tr\bigl(X_0\{X_1,X_2\}\bigr),&w\text{ odd}.
 \end{cases}
 \label{eq:all-order-cubic-shape}
\end{equation}
Here $\{X_1,X_2\}=X_1X_2+X_2X_1$ is the matrix anticommutator.
The parity choices make the entire cochain alternating.  There is
no invariant two-cochain of positive momentum degree, since two
conserved momenta are $p,-p$ and a single vector has no
nonconstant $SL_2$-invariant polynomial.  Together with
$H^3_{\deg=2w}=0$, this proves that the differential is injective
on the invariant cubic space for $w\ge2$.

For specialization, choose $a/N$ to be either fixed square root in
\eqref{eq:all-order-leading-root}.  The inverse trace pairing on
traceless matrices has the exact completeness identity
\begin{equation}
 \sum_b\tr(XT_b)\tr(T^bY)
   =\tr(XY)-\frac1N\tr(X)\tr(Y),
 \label{eq:all-order-completeness}
\end{equation}
where $T_b,T^b$ are trace-dual bases of $\fsl_N$.
Each graph consequently has coefficients in
$\mathbb C[N,N^{-1}]$, its graph periods being scalar constants.
The formal trace differential has rational coefficients independent
of $N$.  At each fixed momentum degree, solving on its image
therefore introduces only nonzero rational denominators, not new
polynomials in $N$ in denominators.  Induction constructs the
entire local identity over $\mathbb C[N,N^{-1}]$.
Evaluation on traceless matrices at a positive integer $N$ commutes
with the classical differential, tensor contractions, and graph
composition.  Hence every specialized identity holds, whether or
not additional trace relations exist at that rank.

For $N=2$, the identity
$X_1X_2+X_2X_1=\tr(X_1X_2)\mathbf1_2$ on traceless matrices
annihilates every odd-order cubic in
\eqref{eq:all-order-cubic-shape}.  In particular $J_1=0$.
The even-order specialized counterterms still satisfy the full
local equations.  No inference about the intrinsic cohomology
of the small-rank complex is used.
\end{proof}

\begin{remark}[What is and is not classified]
The uniqueness statement fixes the heat-kernel scheme, the leading
root, and the full-wedge equivariant cubic representatives.  It is
not a statement that arbitrary local field redefinitions or arbitrary
renormalization schemes are identical.  The formal series need not
converge analytically in $\hbar$.  Its coefficients are local and of
finite derivative order at each order, which is the perturbative
meaning of quantization used here.
\end{remark}

\subsection{A finite evaluation determining every higher cubic}

The all-order theorem also gives an exact extraction formula for the
counterterm coefficient.  This is useful independently of any
choice of basis for four-cochains.  For $N\ge4$, evaluate an
order-$w$ equation on the color matrices
\begin{equation}
  (X_0,X_1,X_2,X_3)=(E_{12},E_{23},E_{34},E_{41})
\end{equation}
and on the conserved momenta
\begin{equation}
  (p_0,p_1,p_2,p_3)=((1,0),(0,1),(1,1),(-2,-2)).
\end{equation}
Here $E_{ab}$ is the matrix with its only nonzero entry equal to one
in row $a$ and column $b$.  The differential of the unscaled cubic
in \eqref{eq:all-order-cubic-shape} evaluates to
\begin{align}
 -[p_2p_3]^w+[p_1p_2]^w+[p_3p_0]^w-[p_0p_1]^w
   &=2^w+(-1)^w-1\notag\\
   &=\begin{cases}2^w,&w\text{ even},\\
                    2^w-2,&w\text{ odd}.
      \end{cases}
 \label{eq:all-order-extraction-denominator}
\end{align}
It is nonzero for every $w\ge2$.  Dividing the evaluated right-hand
side of the cochain equation $d i_w=\text{known lower-order terms}$
by this integer determines its unique invariant cubic primitive.
This extraction does not prove closedness; the master-equation
identity and the cohomology calculation supply that logically prior
step.  Explicit four-loop periods and the existing independent
matrix calculations are recorded in~\cite{RepoHigher}.

\section{The complete fourth-order coefficient}
\label{sec:fourth-order}

We now fix the Wick convention $\epsilon U=1$. The factor
$(\epsilon\pi\nu)^2=1/(4096\pi^8)$ will be left explicit.
At order $\hbar^4$ there are tree graphs, marked two-loop
theta graphs, and bare four-loop graphs. Odd geometric loop
numbers vanish, irrespective of the counterterm orders at
the vertices.

\subsection{Trees and marked theta graphs}

Use the four matrix units and momenta of
\eqref{eq:all-order-extraction-denominator}. The differential
of the unscaled fourth-order cubic evaluates to $16$.
On these inputs the two tree contractions are
\[
 \{D,[p_0p_1]^3\tr(X_0\{X_1,X_2\})\}=10,\qquad
 \frac12\{[p_0p_1]^2f,[p_0p_1]^2f\}=-4.
\]
For the marked theta graphs the following entries include
the sum over allowed marked vertices and the graph
automorphism factors, with the displayed vertex coupling
and the common factor $\pi\nu$ removed:
\begin{center}
\renewcommand{\arraystretch}{1.7}
\begin{tabular}{@{}lcc@{}}
\toprule
Theta graph & One $[p_0p_1]^2f$ vertex & Two $D$ vertices\\
\midrule
$(1,3,3)$ & $-\dfrac{499N^2}{2835}-\dfrac{256}{105}$
  & $\dfrac{431N^2}{11340}+\dfrac{718}{2835}$\\
$(2,2,3)$ & $-\dfrac{34}{27}$ & $\dfrac5{63}$\\
\bottomrule
\end{tabular}
\end{center}
Every contraction uses the complete traceless inverse pairing,
including its $-1/N$ term. Substituting
\eqref{eq:first-two-counterterms} and \eqref{eq:i3}, adding the
trees, and dividing by $16$ gives the following contribution
to the primitive coefficient:
\begin{equation}
 (\epsilon\pi\nu)^2
       \frac{21221N^4+85304N^2+192480}{3628800}.
 \label{eq:four-tree-marked}
\end{equation}

\subsection{The fifty bare graphs}

A tight cubic four-loop graph has ten vertices, thirteen
internal edges, and four external fields. The internal degree
of every external vertex is two: an internal degree-one vertex
would violate the tight bound on the remaining induced
subgraph. A cut vertex is also excluded by applying the tight
bound to the two subgraphs that meet there. Suppressing the
four bivalent vertices therefore leaves a biconnected cubic
multigraph on six vertices.

A loop at a core vertex is incompatible with this
biconnectedness. Three parallel edges would isolate two
vertices. The remaining cores are enumerated by their simple
graphs and the matchings that specify doubled edges.
Distributing the four subdivisions, dividing by the core
automorphisms, and checking the tight inequality on every
induced subset leaves exactly fifty isomorphism classes.
The four surviving core types contain $3,8,28,11$ classes.
The oriented edges, subdivision paths, and full automorphism
orders are recorded individually in the accompanying
\texttt{four-loop-graphs.json} data of \cite{RepoHigher}.

The integrated kernels are homogeneous quartics in the three
momentum determinants. On the fixed cyclic matrix-unit input,
the four core sums, divided by $\pi^2$ before the common heat
normalization, have the following coefficients:
\begin{center}
\small
\renewcommand{\arraystretch}{2.1}
\setlength{\tabcolsep}{4pt}
\begin{tabular}{@{}lccc@{}}
\toprule
Graph indices & Coefficient of $1$ & Coefficient of $N^2$
  & Coefficient of $N^4$\\
\midrule
$0$--$2$ & $0$ & $0$ & $\frac1{10500}$\\
$3$--$10$
 & $\frac{101}{945}-\frac{1241\pi^2}{45360}$
 & $\frac{73903}{1020600}-\frac{37\pi^2}{3402}$
 & $\frac{34259}{20412000}-\frac{11\pi^2}{829440}$\\
$11$--$38$
 & $\frac{9796575\pi^2-96791482}{12247200}$
 & $\frac{88921725\pi^2-884922356}{734832000}$
 & $\frac{144949856-14698125\pi^2}{653184000}$\\
$39$--$49$
 & $\frac{93253402-9461505\pi^2}{12247200}$
 & $\frac{734213669-74453325\pi^2}{183708000}$
 & $0$\\
\bottomrule
\end{tabular}
\end{center}
Their exact sum is
\begin{equation}
 -\frac{172}{945}
 +N^2\left(\frac{139229}{48600}-\frac{35701\pi^2}{120960}\right)
 +N^4\left(\frac{760981}{3402000}
                       -\frac{130727\pi^2}{5806080}\right).
 \label{eq:four-bare-sum}
\end{equation}
In this convention the common four-loop heat factor in the
cochain equation is
\[
 \frac1{(8\pi^{5/2})^4}
       =\frac{(\epsilon\pi\nu)^2}{\pi^2}.
\]
The contribution of \eqref{eq:four-bare-sum} to the primitive
is therefore $(\epsilon\pi\nu)^2/16$ times that expression.

\begin{theorem}[Fourth-order counterterm]
\label{thm:i4}
In the stated heat-kernel scheme and Wick convention, the
fourth-order Taylor cochain is
\begin{equation}
\boxed{\begin{aligned}
 i_4={}&\frac{1}{4096\pi^8}
 \left[
   \frac1{24}
   +N^2\left(\frac{1102559}{5443200}
                       -\frac{35701\pi^2}{1935360}\right)\right.\\[-1mm]
 &\left.\hspace{23mm}
   +N^4\left(\frac{4216}{212625}
                       -\frac{130727\pi^2}{92897280}\right)
 \right][p_0p_1]^4\tr(X_0[X_1,X_2]).
\end{aligned}}
\label{eq:i4}
\end{equation}
The entire local order-four equation holds for arbitrary
traceless colors and conserved momenta, and specializes to
every $N\ge2$. The same $i_4$ applies to both choices of $a$.
\end{theorem}

\begin{proof}
Add \eqref{eq:four-tree-marked} to one sixteenth of
\eqref{eq:four-bare-sum}. The constant part is
\[
 \frac{192480}{3628800}-\frac{172}{16\cdot945}
       =\frac1{24};
\]
the $N^2$ and $N^4$ parts give exactly the two displayed
parentheses. Theorem~\ref{thm:all-orders-traceless} proves
closedness and uniqueness before this evaluation is used.
Thus the single evaluation determines the coefficient.
There is also a separate full-cochain check: after reconstructing
the color contractions and all differentiated tree and theta
terms, the rational part and the $\pi^2$ part each agree with
$di_4$ in all $300$ nonzero trace and momentum coordinates,
including their entire Laurent polynomials in $N$.
\end{proof}

\begin{corollary}
The coefficient of
$(\epsilon\pi\nu)^2[p_0p_1]^4f$ in $i_4$ is positive and
transcendental for every integer $N\ge2$.
\end{corollary}
\begin{proof}
The strict inequality $\pi^2<10$ gives
\begin{align*}
 \frac{1102559}{5443200}
          -\frac{35701\pi^2}{1935360}
    &>\frac{787747}{43545600}>0,\\
 \frac{4216}{212625}
          -\frac{130727\pi^2}{92897280}
    &>\frac{2864617}{497664000}>0.
\end{align*}
At each positive integer $N$, the coefficient is a rational
number minus
\[
 \left(\frac{35701N^2}{1935360}
              +\frac{130727N^4}{92897280}\right)\pi^2 ,
\]
whose rational multiplier of $\pi^2$ is nonzero.
Transcendence of $\pi$ then proves the assertion.
\end{proof}

The coefficient is attached to this renormalization scheme
and this choice of local representative. Its transcendence
does not identify a scheme-independent observable.
The full integrated kernels are inherited exact input;
the present verification independently reconstructs their
complete cochain assembly, rather than repeating every
four-loop master integration.

\section{Cancellation by one additional central field}
\label{sec:central-extension}

The hypothesis that the gauge algebra is semisimple is essential to the
no-go statement.  A central field supplies a symmetric invariant cubic
even when the simple algebra has none.  The resulting cancellation must
also satisfy the equations with central external fields.  We compute all
of these equations below.

Extend the simple algebra $\g$ by a central generator $\mathbf 1$ and
extend its invariant form by
\begin{equation}
 \kappa(X,\mathbf1)=0,\qquad
 \kappa(\mathbf1,\mathbf1)=\ell\ne0,\qquad X\in\g.
 \label{eq:central-pairing}
\end{equation}
The symmetric invariant cubic $T$ is specified by
\begin{equation}
\begin{split}
 T(X,Y,\mathbf1)&=\chi\kappa(X,Y),\qquad
 T(\mathbf1,\mathbf1,\mathbf1)=\gamma,\\
 T(X,Y,Z)&=0,\qquad T(X,\mathbf1,\mathbf1)=0,
 \qquad X,Y,Z\in\g.
\end{split}
\label{eq:central-cubic}
\end{equation}
All evaluations not displayed are fixed by symmetry and multilinearity.
Its Taylor cochain is
\begin{equation}
 j(X_0e^{p_0z},X_1e^{p_1z},X_2e^{p_2z})
   =[p_0p_1]T(X_0,X_1,X_2),\qquad p_0+p_1+p_2=0.
 \label{eq:central-j}
\end{equation}

\begin{proposition}[The complete central-field square]
\label{prop:central-square}
The cochain \eqref{eq:central-j} is classically closed.
For four external momenta with sum zero, define the already occurring
metric four-cochain by
\begin{equation}
 G_{kk}=\alt\bigl([p_0p_1][p_2p_3]
                 \kappa(X_0,X_1)\kappa(X_2,X_3)\bigr),
 \qquad X_i\in\g.
\end{equation}
The values of $\tfrac12\{j,j\}$ on every central-input sector are
\begin{equation}
\begin{array}{c|l}
 \text{number of central inputs}&\tfrac12\{j,j\}\\ \hline
 0&\displaystyle \frac{\chi^2}{8\ell}\,G_{kk}\\[2mm]
 1&0\\[1mm]
 2&\displaystyle
 \chi\left(\frac\gamma\ell-\chi\right)
 [p_0p_1][p_2p_3]\kappa(X_0,X_1),
 \quad (X_2,X_3)=(\mathbf1,\mathbf1)\\[2mm]
 3&0\\[1mm]
 4&0.
\end{array}
\label{eq:central-square-sectors}
\end{equation}
The two-central value in any other order follows by alternation.
For $\chi\ne0$, eliminating the two-central obstruction requires
\begin{equation}
 \gamma=\ell\chi.
 \label{eq:central-necessary}
\end{equation}
\end{proposition}

\begin{proof}
The only potentially nonzero classical differential of the mixed cubic
has three simple inputs and one central input.  Its value is
\begin{equation}
 -\chi\kappa([X_0,X_1],X_2)
 \bigl([p_0p_3]+[p_1p_3]+[p_2p_3]\bigr)=0.
\end{equation}
The pure central cubic has zero classical differential.  Thus $dj=0$.

The six $(2,2)$ shuffles in the bracket occur in equal pairs.  After the
factor $1/2$, their sum is the three pair partitions, with signs
$+,-,+$:
\begin{align}
 \tfrac12\{j,j\}
 &= [p_0p_1][p_2p_3]
       \kappa^{-1}\bigl(T(X_0,X_1,\cdot),T(\cdot,X_2,X_3)\bigr)\notag\\
 &\quad-[p_0p_2][p_1p_3]
       \kappa^{-1}\bigl(T(X_0,X_2,\cdot),T(\cdot,X_1,X_3)\bigr)\notag\\
 &\quad+[p_0p_3][p_1p_2]
       \kappa^{-1}\bigl(T(X_0,X_3,\cdot),T(\cdot,X_1,X_2)\bigr).
 \label{eq:central-three-pairings}
\end{align}
For four simple inputs the contracted field is central, so its inverse
pairing is $\ell^{-1}$.  The full alternation defining $G_{kk}$ is eight
times the three pairings in \eqref{eq:central-three-pairings}.  This gives
the first row of \eqref{eq:central-square-sectors}.

For $(X_0,X_1,\mathbf1,\mathbf1)$, the first pairing contributes
$\chi\gamma\kappa(X_0,X_1)/\ell$.  The other pairings each contribute
$\chi^2\kappa(X_0,X_1)$.  The exact determinant identity
\begin{equation}
 [p_0p_1][p_2p_3]-[p_0p_2][p_1p_3]
       +[p_0p_3][p_1p_2]=0
 \label{eq:central-plucker}
\end{equation}
then gives the two-central row.  The same identity kills the
four-central row, whose common color coefficient is $\gamma^2/\ell$.
The one- and three-central rows vanish from
\eqref{eq:central-cubic} before the momentum sum.

To prove necessity of \eqref{eq:central-necessary}, take
$X_0=X_1=X$ with $\kappa(X,X)\ne0$.  These two inputs and the two
central inputs commute.  Every classical boundary vanishes on them,
whereas $[p_0p_1][p_2p_3]$ is not the zero polynomial on conserved
momenta.  Bare diagrams also vanish when any external field is central.
Consequently a nonzero two-central row cannot be removed by an
order-$\hbar^2$ counterterm.
\end{proof}

\begin{theorem}[A full two-loop central cancellation]
\label{thm:central-cancellation}
Use the Killing form on $\g$.  Suppose the original graph tensors obey
the cochain identities
\begin{equation}
 F_{133}+F_{223}=\tfrac12\beta,\qquad
 F_{223}=x\beta+aG_{kk},\qquad
 \beta=d\bigl([p_0p_1]^2\kappa(X_0,[X_1,X_2])\bigr).
 \label{eq:central-graph-assumption}
\end{equation}
Then the theory for $\g\oplus\mathbb C\mathbf1$ has a solution of the
complete local quantum master equation through order $\hbar^2$ with
\begin{equation}
 \boxed{\quad
 \gamma=\ell\chi,\qquad
 \chi^2=\frac{\epsilon U\ell a}{384\pi^4},\qquad
 i_2=\frac{\epsilon U}{\pi^4}
       \left(\frac1{10240}-\frac{x}{3072}\right)
       [p_0p_1]^2\kappa(X_0,[X_1,X_2]).\quad}
 \label{eq:central-solution}
\end{equation}
Here $i_2$ vanishes if one of its arguments is central, and the
coefficients are taken over $\mathbb C$.  The signs $\epsilon$ and $U$
are the tree and loop signs of the main normalization.
\end{theorem}

\begin{proof}
A central external field kills every bare cubic graph because its
adjoint action is zero.  The bare two-loop cochain is therefore the
simple-algebra cochain, extended by zero, and
\begin{equation}
 o_2=-\frac{U}{64\pi^4}
       \left[\left(\frac{x}{48}-\frac1{160}\right)\beta
                  +\frac{a}{48}G_{kk}\right].
 \label{eq:central-bare-anomaly}
\end{equation}
The chosen central cubic cancels the two-central row in
\eqref{eq:central-square-sectors}.  The coefficient of $G_{kk}$ in
$\tfrac12\{j,j\}+\epsilon o_2$ is exactly
\begin{equation}
 \frac{\chi^2}{8\ell}-\frac{\epsilon U a}{3072\pi^4}=0.
\end{equation}
The remaining coefficient is the classical differential of the
displayed $i_2$.  Thus
$-di_2+\tfrac12\{j,j\}+\epsilon o_2=0$ on every input, which is the
full order-$\hbar^2$ equation in the stated descent convention.
The order-$\hbar$ equation was proved above.  The local one-loop
contacts containing $j$ vanish by the odd-loop theorem.
\end{proof}

The hypothesis \eqref{eq:central-graph-assumption} is an identity of
entire cochains.  A numerical contraction detecting one nonzero class
does not by itself establish this hypothesis.  The previously computed
identities in \cite{RepoOriginal} supply examples; the following one is
also recomputed directly in the present paper's certificate.

\begin{corollary}[An explicit $\mathfrak{so}(5)\oplus\mathbb C$ solution]
\label{cor:so5-central}
For $\g=\mathfrak{so}(5)$ with its Killing form,
\begin{equation}
 F_{223}=-\beta-\frac1{18}G_{kk}.
\end{equation}
Consequently
\begin{equation}
 \chi^2=-\frac{\epsilon U\ell}{6912\pi^4},\qquad
 \gamma=\ell\chi,\qquad
 i_2=\frac{13\epsilon U}{30720\pi^4}
          [p_0p_1]^2\kappa(X_0,[X_1,X_2])
 \label{eq:so5-central-solution}
\end{equation}
solves the complete local equation through order $\hbar^2$.
At $\ell=1$ and $\epsilon U=1$, the two first-order choices are
\begin{equation}
 \chi=\gamma=\frac{i}{48\sqrt3\,\pi^2}
 \qquad\text{or}\qquad
 \chi=\gamma=-\frac{i}{48\sqrt3\,\pi^2}.
\end{equation}
\end{corollary}

For completeness, the full graph identities already certified
in \cite{RepoOriginal} for the five exceptional simple algebras
give the following additional instances of
Theorem~\ref{thm:central-cancellation}:
\begin{center}
\renewcommand{\arraystretch}{1.6}
\begin{tabular}{@{}lcc@{}}
\toprule
$\g$ & $\chi^2\pi^4/(\epsilon U\ell)$
 & Coefficient of $[p_0p_1]^2f$ in $\pi^4 i_2/(\epsilon U)$\\
\midrule
$G_2$ & $-\dfrac5{36864}$ & $\dfrac{71}{368640}$\\
\addlinespace[3pt]
$F_4$ & $-\dfrac5{124416}$ & $\dfrac{41}{829440}$\\
\addlinespace[3pt]
$E_6$ & $-\dfrac1{36864}$ & $\dfrac{11}{368640}$\\
\addlinespace[3pt]
$E_7$ & $-\dfrac1{62208}$ & $\dfrac{11}{829440}$\\
\addlinespace[3pt]
$E_8$ & $-\dfrac1{115200}$ & $\dfrac1{460800}$\\
\bottomrule
\end{tabular}
\end{center}
In each row $\gamma=\ell\chi$. These are algebraic consequences
of the cited complete cochain identities. The new independent
original-graph replay described next is for $\mathfrak{so}(5)$.

\paragraph{Exact certificate.}
The script \texttt{checks/central\_extension\_exact.py} constructs the
ten matrices $E_{ij}-E_{ji}$ in $\mathfrak{so}(5)$ and checks their
complete commutator table, Jacobi identity, and Killing form
$\kappa=-6\delta$.  It contracts the six original structure tensors of
each theta graph, with no graph-engine imports, then performs all
$24$ alternations and checks all six coefficients of
$[p_0p_1]^2,[p_0p_2]^2,[p_1p_2]^2,
[p_0p_1][p_0p_2],[p_0p_1][p_1p_2],[p_0p_2][p_1p_2]$.
Each graph identity is checked on $60000$ color--momentum entries.
The central extension has respectively $10000,4000,600,40,1$
ordered color assignments with $0,1,2,3,4$ central inputs.
The full bracket certificate checks $351384$ coefficient entries;
the assembled quantum master equation checks $87846$ entries.
The arithmetic is integral after explicit denominator clearing.

This establishes a two-loop cancellation for a reductive algebra
whose simple factor is obstructed on its own.  It does not establish
an all-order quantization for that reductive algebra: higher
obstructions require the corresponding local cohomology and graph
consistency analysis.

\section{Further consequences and remaining questions}

The obstruction is sensitive to the available color tensors.
For a non-type-$A$ simple algebra, the
entire positive-derivative first-order cohomology vanishes,
so changing derivative order or dropping rotational
equivariance cannot repair the two-loop anomaly.
For traceless matrices the symmetric cubic supplies the
necessary first-order square, and the higher invariant
obstruction groups vanish. A central field changes the first
statement by creating an additional invariant cubic; its
mixed external equations fix the relation
$\gamma=\ell\chi$.

Three mathematical questions remain beyond the results
proved here. First, the higher local cohomology and the
marked-graph identities for the interacting central
extensions must be determined before an all-order theorem
can be stated for them. Second, the intrinsic cohomology
at small matrix ranks can contain extra trace relations;
specializing a full identity establishes existence without
classifying those additional relations. Third, the sequence
of higher graph periods remains to be computed even though
the invariant cohomology already proves existence and
uniqueness of each formal counterterm. None of these
questions requires changing the exact two-loop classification.

\appendix
\section{Verification, provenance, and reproduction}
\label{sec:verification}

\subsection{Analytic arguments and exact certificates}

The all-order assertion is proved by the contact theorem,
the cyclic and two-trace cohomology calculations, and the
master-equation induction. It is not inferred from checking
a finite number of perturbative orders. Likewise the
rank-independent no-go result uses the highest-root argument
and the polynomial primitive in
Lemma~\ref{lem:f-channel-all-degrees}.

The reproducible calculations accompanying this version
check the following independent ingredients:
\begin{enumerate}
\item The regulated kernel characters and the two theta
normalizations, with exact rational and gamma-function
arithmetic.
\item The original two-loop graphs, the classical differential,
the contracted brackets, and the marked two-loop terms
through order three on exact traceless $3\times3$ matrices.
The supplied replay also includes deliberately incorrect
relative signs that fail the identities.
\item The complete fourth-order formal trace identity,
using all fifty stored integrated geometric kernels.
Both rational and $\pi^2$ coefficient dictionaries have
zero remainder in all $300$ nonzero coordinates.
\item The Weyl-character cubic multiplicities for $32$
independently generated root systems, including $14$
complete Weyl-group sums. The largest complete sum has
$51840$ elements in type $E_6$.
The polynomial first-order calculation is checked in all
$25$ bidegrees with each entry between zero and four;
the explicit primitive signs are checked on all $27$
ordered color triples of $\fsl_2$.
\item The complete central-extension identities, including
the original $\fso_5$ graph contractions and all central
external sectors, with the coefficient counts stated
after Corollary~\ref{cor:so5-central}.
\item An exhaustive independent enumeration of all loopless
six-vertex cubic multigraphs, all four-subdivision distributions,
and every induced-subset bound, giving the same fifty graphs
and automorphism orders. The exact four-loop core sums,
tree and marked-graph assembly, and the rational
positivity bounds in Theorem~\ref{thm:i4} are also checked.
\item An independent polygon-coordinate construction of the
cyclic rotation and deletion matrices at momentum degrees
$2,3,4,5$. Exact ranks over each of the fields with
$1000003$ and $1000033$ elements give cyclic cohomology
dimensions $(1,0)$ in cochain degrees $(2,3)$ at momentum
degree two and $(0,0)$ at degrees three through five.
These bounded checks support the all-degree analytic argument.
\end{enumerate}
The accompanying \texttt{README.md} gives the commands,
dependencies, and outputs. The \texttt{Makefile} builds
the paper and runs the certificates. The manifest records
the source revision and hashes of the inherited kernel
inputs. The JSON and text outputs are included with the
source and compiled document.

\subsection{Normalization and evidence boundaries}

The numerical coefficients use the explicit orientation,
Taylor-cochain factorials, full alternation, and heat
normalization of Sections~\ref{sec:setting} and
\ref{sec:two-loop}. The original companion distinguishes
its relative graph-sign certificates from a complete
end-to-end audit of every Batalin--Vilkovisky action
convention. We preserve that distinction by recording
$\epsilon$ and $U$ and by fixing $\epsilon U=1$ when
stating the evaluated fourth-order coefficient.
The relative cochain identities have been replayed;
the classification over $\mathbb C$ does not depend on
the choice of these overall signs.

The fourth-order kernel tables are inherited from
the exact integrations in \cite{RepoHigher}. This
version checks their graph census, exact assembly,
and full color--momentum identity. It does not claim
to have rerun the external integration-by-parts
master-integral reduction software.

The manuscript concerns the commutative classical
theory in the stated space of full exterior products
of holomorphic jets. It proves formal local
perturbative quantization with its specified
symmetries. It makes no assertion about convergence
of the formal series, nonperturbative path
integrals, or global gauge transformations on
compact five-manifolds.


\begin{thebibliography}{99}
\raggedright
\bibitem{WW}
M.~Wang and B.~R.~Williams,
\emph{On the renormalization and quantization of
topological--holomorphic field theories},
\href{https://arxiv.org/abs/2407.08667v3}{arXiv:2407.08667v3};
Propositions 3.18, 3.21--3.23.
\bibitem{CostelloBV}
K.~Costello,
\emph{Renormalisation and the Batalin--Vilkovisky formalism},
\href{https://arxiv.org/abs/0706.1533}{arXiv:0706.1533};
Sections 12--13.
\bibitem{CostelloOmega}
K.~Costello,
\emph{M-theory in the Omega-background and 5-dimensional
non-commutative gauge theory},
\href{https://arxiv.org/abs/1610.04144}{arXiv:1610.04144};
Sections 5 and 9.
\bibitem{UrichukWalton}
A.~Urichuk and M.~A.~Walton,
\emph{Adjoint affine fusion and tadpoles},
\href{https://arxiv.org/abs/1606.03842}{arXiv:1606.03842};
equation (18).
\bibitem{MSV}
R.~L.~Mkrtchyan, A.~N.~Sergeev and A.~P.~Veselov,
\emph{Casimir eigenvalues for universal Lie algebra},
\href{https://arxiv.org/abs/1105.0115}{arXiv:1105.0115}.
\bibitem{LLC}
Li, Liu and Chang,
\emph{Holomorphic Feynman diagrams in four dimensions and the
quantum BPS ring}, source manuscript
\href{https://github.com/lbotao399-stack/holomorphic-topological-feynman-diagrams/blob/b12b46ef484a1b002ddf92a8d776d44bab0a4f5c/sources/LLC_main.tex}{\texttt{sources/LLC\_main.tex}}
in the computational repository.
\bibitem{RepoOriginal}
\emph{Commutative Chern--Simons theory on
$\mathbb R\times\mathbb C^2$ through order $\hbar^3$
and the two-loop theta kernels on
$\mathbb R^T\times\mathbb C^H$},
\href{https://github.com/lbotao399-stack/holomorphic-topological-feynman-diagrams/blob/b12b46ef484a1b002ddf92a8d776d44bab0a4f5c/notes/flagship/flagship.tex}{\texttt{notes/flagship/flagship.tex}},
with its computational certificates and referee checks.
\bibitem{RepoHigher}
\emph{Higher orders of five-dimensional commutative
Chern--Simons theory},
\href{https://github.com/lbotao399-stack/holomorphic-topological-feynman-diagrams/blob/b12b46ef484a1b002ddf92a8d776d44bab0a4f5c/deepen/o3/higher_orders/README.md}{\texttt{deepen/o3/higher\_orders/README.md}},
including the four-loop reduction and complete cochain replay.
\bibitem{RepoContacts}
\emph{Local contacts and changes of first-order representative},
\href{https://github.com/lbotao399-stack/holomorphic-topological-feynman-diagrams/blob/b12b46ef484a1b002ddf92a8d776d44bab0a4f5c/deepen/h1/verify/regulated_contacts.md}{\texttt{deepen/h1/verify/regulated\_contacts.md}}.
\end{thebibliography}
\end{document}
